\appendix

% ============================================================
\chapter{手法選択フローチャート}\label[appendix]{app:A}
% ============================================================
アウトカムの型と観測の構造から手法を選ぶための一覧である．
本書で定義した手法だけを載せている．
まず\cref{fig:A-flow}で候補を絞り，\cref{tab:A-select}で推定対象と対応する節を確認する．

\begin{figure}[!htbp]
\centering
\begin{tikzpicture}[
  typ/.style={draw=bkblue,fill=bkblue!8,rounded corners=2pt,font=\small\bfseries,text width=17mm,align=center,minimum height=10mm},
  lst/.style={draw=bkgray!60,fill=white,rounded corners=2pt,font=\small,text width=132mm,align=left,inner sep=4pt},
  hd/.style={font=\small\bfseries,text=bkblue},
  arr/.style={-{Stealth[length=2.2mm]},bkblue,thick}]
\node[lst,anchor=north west] (cl) at (0,0) {%
\textbf{独立2群}：正規・等分散 → 2標本 $t$，不等分散 → Welch，非正規・外れ値 → 順位和（Mann--Whitney）\\
\textbf{対応あり}：正規 → 対応のある $t$，対称 → 符号付き順位，形の仮定なし → 符号検定\\
\textbf{3群以上}：一元配置分散分析 / Kruskal--Wallis，用量の順序 → 線形対比\\
\textbf{対数正規（PK）・生物学的同等性} → 対数変換して差の 90\%CI};
\node[lst,anchor=north west] (bl) at ([yshift=-5mm]cl.south west) {%
\textbf{独立2群}：$\chi^2$（大標本）/ Fisher（小標本），RD・RR・OR の区間（対数スケールの Wald）\\
\textbf{対応あり}（前後・ペア・同一患者に2検査）→ McNemar，条件付き OR $n_{12}/n_{21}$\\
\textbf{層別・交絡} → 層別 Cochran（CMH）検定，OR の均一性 → Woolf の検定，\textbf{多数の共変量} → ロジスティック回帰\\
\textbf{順序のある $K$ 群} → Cochran--Armitage，\textbf{診断} → 感度・特異度・PPV・NPV・ROC/AUC};
\node[lst,anchor=north west] (pl) at ([yshift=-5mm]bl.south west) {%
\textbf{人時間つき件数} → ポアソン率 $\hat\lambda=Y/T$，$\V[\log\hat\lambda]\approx1/Y$\\
\textbf{小標本の割合＋事前情報} → Beta--Binomial の事後分布・事後確率};
\node[lst,anchor=north west] (sl) at ([yshift=-5mm]pl.south west) {%
\textbf{記述} → Kaplan--Meier（分散は Greenwood の近似公式）\\
\textbf{2群比較} → ログランク検定\\
\textbf{調整・効果の大きさ} → Cox（HR，部分尤度），ハザード一定 → 指数モデル\\
\textbf{競合イベント} → 原因別ハザードと累積発生関数で解釈（原因別 Cox と Fine--Gray の係数の違い）};
\node[lst,anchor=north west] (ml) at ([yshift=-5mm]sl.south west) {%
\textbf{連続・反復測定} → 線形混合モデル（変量切片＝複合対称），共分散構造は AIC で選択\\
\textbf{欠測} → MCAR なら完全ケース，MAR なら尤度法（EM アルゴリズム）};
\foreach \l/\t/\n in {cl/連続/c,bl/2 値/b,pl/{カウント\\・率}/p,sl/生存時間/s,ml/{反復測定\\・欠測}/m}{
  \node[typ,left=6mm of \l] (\n) {\t};
  \draw[arr] (\n.east)--(\l.west);}
\node[hd,anchor=south] at ([yshift=2mm]c.north) {アウトカム};
\node[hd,anchor=south west] at ([yshift=2mm]cl.north west) {状況 → 手法};
\end{tikzpicture}
\caption{アウトカムからの手法選択}\label{fig:A-flow}
\end{figure}
\clearpage

\begin{center}\small
\begin{longtable}{@{}L{36mm}L{30mm}L{52mm}c@{}}
\caption{データ構造・推定対象・手法の対応}\label{tab:A-select}\\
\toprule
データ構造 & 推定対象 & 手法 & 節\\
\midrule
\endfirsthead
\toprule
データ構造 & 推定対象 & 手法 & 節\\
\midrule
\endhead
\bottomrule
\endfoot
連続・対応あり & 差の平均・中央値 & 対応のある $t$，符号付き順位，符号 & \ref{sec:03-signrank}\\
連続・独立2群 & 平均差，$\Prob(X>Y)$ & 2標本 $t$，Welch，順位和（Mann--Whitney） & \ref{sec:03-ranksum}\\
連続・3群以上 & 群平均の差 & 一元配置分散分析，Kruskal--Wallis & \ref{sec:14-anova}，\ref{subsec:03-kw}\\
連続・多時点 & 時点別平均差・変化量 & 線形混合モデル，共分散構造の AIC & \ref{sec:08-lmm}\\
2 値・独立2群 & RD，RR，OR & $\chi^2$，Fisher，対数スケールの Wald 区間 & \ref{subsec:04-var}\\
2 値・$K$ 群 & 一様性・傾向 & $\chi^2_{K-1}$，Cochran--Armitage & \ref{sec:04-trend}\\
2 値・対応あり & 周辺割合の差，条件付き OR & McNemar，LRT，$n_{12}/n_{21}$ & \ref{sec:04-mcnemar}\\
2 値・層別 & 全層で関連なしの検定 & 層別 Cochran（CMH）検定 & \ref{sec:04-mh}\\
2 値・層別 & OR の均一性 & Woolf の検定 & \ref{subsec:07-homog}\\
2 値・共変量 & 調整 OR，予測確率 & ロジスティック回帰，LRT，AIC & \ref{ch:06}\\
診断検査 & 感度・特異度・PPV・AUC & ROC，McNemar，多項分布＋デルタ法 & \ref{ch:05}\\
件数＋人時間 & 率 & ポアソン率の MLE & \ref{sec:17-poisson}\\
生存時間 & $S(t)$，中央値 & KM，Greenwood（近似） & \ref{sec:09-km}\\
生存時間・2群 & 分布の差 & ログランク & \ref{sec:09-logrank}\\
生存時間・共変量 & HR & Cox，部分尤度 & \ref{sec:10-cox}\\
生存時間・ハザード一定 & ハザード・HR & 指数モデル，デルタ法 & \ref{sec:10-exp}\\
競合リスク & 原因別ハザード，累積発生関数 & 原因別 Cox と Fine--Gray の解釈 & \ref{sec:10-compete}\\
同等性・生物学的同等性 & 差・比がマージン内 & $(1-2\alpha)$CI による判定 & \ref{sec:13-be}\\
中間解析 & 早期中止 & $\alpha$ 消費関数 & \ref{sec:14-gsd}\\
多重比較・多重評価項目 & FWER 制御 & Bonferroni，Holm，閉検定手順 & \ref{sec:14-mult}\\
用量反応 & 傾向 & 線形対比 & \ref{sec:14-contrast}\\
欠測あり & 完全データでの効果 & 完全ケース（MCAR），尤度法・EM（MAR） & \ref{sec:16-methods}\\
小標本＋事前情報 & 事後確率・予測 & Beta--Binomial & \ref{sec:17-betabin}\\
\end{longtable}
\end{center}

% ============================================================
\chapter{必須公式集}\label[appendix]{app:B}
% ============================================================
各公式の右欄は「いつ使うか」である．記号は本文の定義に従う．
誘導付きでしか出題されていない概念の公式は載せていない（\cref{app:F}）．

\begin{center}\small
\renewcommand{\arraystretch}{1.5}
\begin{longtable}{@{}L{25mm}L{80mm}L{45mm}@{}}
\toprule
項目 & 公式 & いつ使うか\\
\midrule
\endfirsthead
\toprule
項目 & 公式 & いつ使うか\\
\midrule
\endhead
\bottomrule
\endfoot
\multicolumn{3}{@{}l}{\textbf{効果指標・分割表}}\\
RD & $\hat\pi_1-\hat\pi_0$，$\V=\frac{\hat\pi_1(1-\hat\pi_1)}{n_1}+\frac{\hat\pi_0(1-\hat\pi_0)}{n_0}$ & 絶対的な効果．NNT $=1/\mathrm{ARR}$（利益），NNH $=1/\mathrm{ARI}$（害）\\
RR & $\frac{a/n_1}{c/n_0}$，$\V[\log\widehat{\RR}]\approx\frac1a-\frac1{n_1}+\frac1c-\frac1{n_0}$ & コホート・RCTの相対効果\\
OR & $\frac{ad}{bc}$，$\V[\log\widehat{\OR}]\approx\frac1a+\frac1b+\frac1c+\frac1d$ & 症例対照研究，ロジスティック\\
区間 & $\exp\{\log\hat\theta\pm z_{\alpha/2}\widehat{\SE}\}$（対数スケールの Wald 区間を逆変換） & RR・OR の信頼区間\\
OR と RR & $\OR=\RR\frac{1-\pi_0}{1-\pi_1}$ & 稀な疾患なら OR $\approx$ RR\\
Pearson $\chi^2$ & $\frac{N(ad-bc)^2}{n_1n_0m_1m_0}\approx\chi^2_1$ & 独立2群の比率の比較\\
一様性 & $\frac{\sum n_i(p_i-p)^2}{p(1-p)}\approx\chi^2_{K-1}$ & $K$ 群の比率\\
傾向性 & $\frac{\{\sum y_i(x_i-\bar x)\}^2}{p(1-p)\sum n_i(x_i-\bar x)^2}\approx\chi^2_1$ & 用量に順序がある比率\\
McNemar & $\frac{(n_{12}-n_{21})^2}{n_{12}+n_{21}}\approx\chi^2_1$，$\V[\hat\delta]=\frac{\pi_{12}+\pi_{21}-(\pi_{12}-\pi_{21})^2}{n}$ & 対応のある2 値\\
条件付き OR & $n_{12}/n_{21}$（不一致ペアの比） & マッチドペア\\
超幾何 & $\E[a]=\frac{n_1m_1}{N}$，$\V[a]=\frac{n_1n_0m_1m_0}{N^2(N-1)}$ & CMH 検定，ログランク\\
CMH & $\chi^2_{\mathrm{CMH}}=\frac{\{\sum_k(a_k-\E[a_k])\}^2}{\sum_k\V[a_k]}\approx\chi^2_1$ & 層別表で関連なしの検定\\
Woolf & $Q=\sum_kw_k(\hat\theta_k-\hat\theta_w)^2\approx\chi^2_{K-1}$，$w_k=1/v_k$ & 層間の OR の均一性\\
\multicolumn{3}{@{}l}{\textbf{診断}}\\
感度・特異度 & $\Prob(+\mid D)$，$\Prob(-\mid\bar D)$ & 検査性能（有病率に依存しない）\\
PPV・NPV & $\frac{Se\,p}{Se\,p+(1-Sp)(1-p)}$，$\frac{Sp(1-p)}{Sp(1-p)+(1-Se)p}$ & 有病率 $p$ の集団での的中率\\
AUC & $\Prob(X_1>X_0)$，経験 $\frac{1}{n_1n_0}\sum\sum I(x_{1i}>x_{0j})$，2正規 $\Phi\bigl(\frac{\mu_1-\mu_0}{\sqrt{\sigma_1^2+\sigma_0^2}}\bigr)$ & ROC 曲線の要約\\
対応のある感度 & $\V[\hat\delta]=\frac1n\{\pi_{+-}+\pi_{-+}-(\pi_{+-}-\pi_{-+})^2\}$ & 同一患者に2検査\\
\multicolumn{3}{@{}l}{\textbf{回帰・混合モデル}}\\
ロジスティック & $\ell=\sum\{y_i\bm x_i\transp\bm\beta-\log(1+e^{\bm x_i\transp\bm\beta})\}$，$\bm U=\bm X\transp(\bm y-\bm\pi)$，$\bm I=\bm X\transp\bm W\bm X$ & 2 値アウトカムの調整 OR\\
LRT & $2(\ell_1-\ell_0)\approx\chi^2_{\text{パラメータ数の差}}$ & 入れ子モデルの比較\\
AIC & $-2\ell(\hat{\bm\theta})+2k$ & モデル・共分散構造の選択\\
混合モデル & $\bm V=\bm Z\bm G\bm Z\transp+\bm R$，$\hat{\bm\beta}=(\bm X\transp\bm V^{-1}\bm X)^{-1}\bm X\transp\bm V^{-1}\bm Y$ & 反復測定\\
\multicolumn{3}{@{}l}{\textbf{順位検定}}\\
符号 & $T\sim\Bin(n,\frac12)$，$\E=\frac n2$，$\V=\frac n4$ & 対応あり・分布の仮定なし\\
符号付き順位 & $\E[W^+]=\frac{n(n+1)}4$，$\V=\frac{n(n+1)(2n+1)}{24}$ & 対応あり・対称分布\\
順位和 & $\E[W]=\frac{m(N+1)}2$，$\V=\frac{mn(N+1)}{12}$，$W=U+\frac{m(m+1)}2$，$\E_0[U]=\frac{mn}2$ & 独立2群・非正規\\
Kruskal--Wallis & $H=\frac{12}{N(N+1)}\sum n_i\bigl(\bar R_i-\frac{N+1}2\bigr)^2\approx\chi^2_{k-1}$ & 3群以上・非正規\\
\multicolumn{3}{@{}l}{\textbf{生存時間}}\\
基本関係 & $h=\frac fS$，$H=-\log S$，$S=e^{-H}$，$f=hS$ & 分布の変換\\
KM & $\hat S(t)=\prod_{t_{(j)}\le t}\bigl(1-\frac{d_j}{n_j}\bigr)$ & 打ち切りのある生存曲線\\
Greenwood & $\widehat\V[\hat S]\approx\hat S^2\sum\frac{d_j}{n_j(n_j-d_j)}$（近似公式） & KM の分散\\
ログランク & $E_{1j}=\frac{d_jn_{1j}}{n_j}$，$V_j=\frac{d_jn_{1j}n_{0j}(n_j-d_j)}{n_j^2(n_j-1)}$，$\frac{(O-E)^2}{V}\approx\chi^2_1$ & 生存曲線の2群比較\\
Cox & $h_0(t)e^{\bm\beta\transp\bm x}$，$L=\prod_j\frac{e^{\bm\beta\transp\bm x_{(j)}}}{\sum_{R_j}e^{\bm\beta\transp\bm x_l}}$，$S=S_0^{\exp(\bm\beta\transp\bm x)}$ & HR の推定・調整\\
指数モデル & $L=\lambda^de^{-\lambda W}$，$\hat\lambda=\frac dW$，$I=\frac d{\lambda^2}$，$\V[\log\hat\lambda]\approx\frac1d$ & ハザード一定・率\\
競合リスク & $S=\exp\{-\sum_j\Lambda_j\}$，$F_j=\int_0^t\lambda_jS$，$F_j(t)\le1-e^{-\Lambda_j(t)}$ & $1-$KM の過大推定\\
\multicolumn{3}{@{}l}{\textbf{設計・臨床試験}}\\
検出力 & $\mu=z_\alpha+z_\beta$，検出力 $\Phi(\mu-z_\alpha)$ & すべての症例数計算の基本\\
平均の差 & $n=\frac{2\sigma^2(z_{\alpha/2}+z_\beta)^2}{\Delta^2}$（各群） & 連続アウトカム\\
比率の差 & $n=\frac{\{z_{\alpha/2}\sqrt{2\bar\pi\bar q}+z_\beta\sqrt{\pi_1q_1+\pi_0q_0}\}^2}{(\pi_1-\pi_0)^2}$ & 2 値アウトカム\\
イベント数 & 指数：$r=\frac{2(z_{\alpha/2}+z_\beta)^2}{(\log\HR)^2}$（各群），ログランク：$D=\frac{(z_{\alpha/2}+z_\beta)^2}{p(1-p)(\log\HR)^2}$ & 生存時間試験\\
同等性 & $(1-2\alpha)$CI $\subset(-\Delta,\Delta)$，BE：GMR の90\%CI $\subset(0.80,1.25)$ & 同等性・生物学的同等性\\
$\alpha$ 消費 & $\Corr[Z(t_1),Z(t_2)]=\sqrt{t_1/t_2}$，増分 $\alpha(t_k)-\alpha(t_{k-1})$ & 中間解析\\
多重性 & Bonferroni $\alpha/m$，Holm $p_{(k)}\le\alpha/(m-k+1)$ & 複数の検定\\
分散分析 & $F=\frac{S_A/(a-1)}{S_E/(N-a)}\sim F(a-1,N-a)$ & 3群以上の平均\\
対比 & $Z_c=\frac{\sum c_j\bar Y_j}{\sqrt{\sigma^2\sum c_j^2/n_j}}$，直線傾向には $c_j\propto x_j-\bar x$ & 用量反応\\
\multicolumn{3}{@{}l}{\textbf{欠測・ポアソン・ベイズ}}\\
EM（ABO） & E：$\E[f_{AA}\mid n_A]=n_A\frac{p}{p+2r}$，M：遺伝子の数え上げ & 不完全データの MLE\\
ポアソン & $\hat\lambda=\frac YT$，$\V[\log\hat\lambda]\approx\frac1Y$ & 人時間つき件数\\
Beta--Binomial & $\Be(a+y,b+n-y)$，事後平均 $\frac{a+y}{a+b+n}$ & 割合の Bayes 推測\\
\end{longtable}
\end{center}

% ============================================================
\chapter{医薬生物学用語集}\label[appendix]{app:C}
% ============================================================
\begin{center}\small\renewcommand{\arraystretch}{1.05}
\begin{longtable}{@{}L{32mm}L{40mm}L{78mm}@{}}
\toprule
用語 & 英語 & 説明\\
\midrule
\endfirsthead
\toprule
用語 & 英語 & 説明\\
\midrule
\endhead
\bottomrule
\endfoot
評価項目 & endpoint & 治療効果を評価する変数．主要評価項目は事前に規定する．複数（co-primary など）とする場合は，成功の条件と多重性の処理も事前に規定する\\
代替評価項目 & surrogate endpoint & 真の評価項目（死亡など）の代わりに用いる検査値など\\
イベント & event & 生存時間解析で観察する事象（死亡・再発）\\
打ち切り & censoring & イベント前に観察が終わること．右側打ち切りが基本\\
独立な打ち切り & independent censoring & 打ち切りがその後の予後と無関係であること\\
リスク集合 & risk set & ある時点の直前にイベントも打ち切りもない人の集合\\
ハザード & hazard & 生存者の瞬間的なイベント発生率 $f/S$\\
累積ハザード & cumulative hazard & $H(t)=\int_0^th=-\log S(t)$\\
ハザード比 & hazard ratio (HR) & 2群のハザードの比．比例ハザードなら時間によらず一定\\
比例ハザード & proportional hazards & HR が時間によらず一定であること\\
中央生存時間 & median survival time & $S(t)=0.5$ となる時点\\
競合リスク & competing risk & あるイベントの発生により他のイベントが起こりえなくなる状況\\
原因別ハザード & cause-specific hazard & 無事な人での原因 $j$ の瞬間的発生率\\
累積発生関数 & cumulative incidence function (CIF) & $\Prob(T\le t,J=j)$\\
部分分布ハザード & subdistribution hazard & CIF に対応するハザード（Fine--Gray）\\
無作為化 & randomization & 治療を確率的に割り付けること．交絡を防ぐ\\
層別無作為化 & stratified randomization & 予後因子の層ごとに割り付けのバランスをとる\\
盲検化 & blinding & 割付を被験者・医師・評価者に知らせないこと\\
二重盲検 & double-blind & 被験者と医師の双方が割付を知らない\\
プラセボ & placebo & 有効成分を含まない偽薬\\
並行群間比較 & parallel group design & 各被験者が1つの治療のみを受けるデザイン\\
クロスオーバー & crossover design & 各被験者が時期を変えて複数の治療を受けるデザイン\\
休薬期間 & washout period & クロスオーバーで持ち越し効果を除くための期間\\
持ち越し効果 & carryover effect & 前の時期の治療効果が次の時期に残ること\\
優越性 & superiority & 対照より優れること\\
同等性 & equivalence & 差がマージン $(-\Delta,\Delta)$ 内にあること\\
非劣性 & non-inferiority & 対照より $\Delta$ を超えて劣らないこと\\
非劣性マージン & non-inferiority margin & 許容できる最大の劣り幅\\
生物学的同等性 & bioequivalence (BE) & 後発品と先発品の吸収の速度・量が同等であること\\
幾何平均比 & geometric mean ratio (GMR) & 対数スケールの平均差を逆変換した比\\
分析感度 & assay sensitivity & 試験が治療間の差を検出できる能力\\
ITT & intention-to-treat & 割り付けられた群のまま全例を解析する原則\\
推定対象 & estimand & 推定したい治療効果の定義（ICH E9(R1)）\\
中間事象 & intercurrent event & 治療開始後に起こり，評価項目の解釈または存在に影響する事象\\
感度 & sensitivity & 罹患者が陽性になる確率\\
特異度 & specificity & 非罹患者が陰性になる確率\\
陽性的中率 & positive predictive value (PPV) & 陽性者が罹患している確率\\
陰性的中率 & negative predictive value (NPV) & 陰性者が罹患していない確率\\
有病率 & prevalence & 集団で疾患を有する割合\\
尤度比 & likelihood ratio (LR) & $Se/(1-Sp)$ など．検査後オッズ＝LR×検査前オッズ\\
ROC 曲線 & receiver operating characteristic curve & (偽陽性率, 真陽性率) の軌跡\\
AUC & area under the curve & ROC 曲線下面積 $=\Prob(X_1>X_0)$．薬物動態では血中濃度曲線下面積\\
リスク & risk & 一定期間の発生確率（累積発生割合）\\
率 & rate & 人時間あたりの発生数\\
オッズ比 & odds ratio (OR) & オッズの比．症例対照研究で推定可能\\
リスク比 & risk ratio (RR) & リスクの比．相対リスク\\
リスク差 & risk difference (RD) & リスクの差\\
NNT & number needed to treat & $1/\mathrm{ARR}$（ARR：絶対リスク減少）．1人の発症を防ぐのに治療が必要な人数\\
NNH & number needed to harm & $1/\mathrm{ARI}$（ARI：絶対リスク増加）．治療により1人余分に害が生じる人数\\
コホート研究 & cohort study & 曝露で分けて追跡する観察研究\\
症例対照研究 & case-control study & 症例と対照を集めて過去の曝露を比べる研究\\
交絡 & confounding & 曝露と結果の両方に関連する因子による偏り\\
効果修飾 & effect modification & 効果の大きさが第3の変数で変わること\\
層別 & stratification & 交絡因子の値ごとに分けて解析すること\\
マッチング & matching & 背景の似た個体を組にすること\\
非可縮性 & non-collapsibility & 交絡がなくても層別と併合で値が変わる性質（OR，HR）\\
Simpson のパラドックス & Simpson's paradox & 併合した関連が層別の関連と逆転する現象\\
欠測 & missing data & 予定された測定が得られないこと\\
MCAR・MAR・MNAR & missing completely at random / at random / not at random & 欠測がデータと無関係／観測値のみに依存／欠測値自体に依存\\
完全ケース解析 & complete-case analysis & 欠測のない個体のみを用いる解析\\
感度分析 & sensitivity analysis & 仮定を変えたときの結論の頑健性の確認\\
中間解析 & interim analysis & 試験途中で行う解析\\
情報時間 & information time & 最終解析に対する情報量の割合\\
$\alpha$ 消費関数 & alpha spending function & 中間解析で使う第1種の過誤の配分を定める関数\\
早期中止 & early stopping & 有効・無効・安全性により試験を途中で終えること\\
多重性 & multiplicity & 複数の検定による第1種の過誤の増大\\
FWER & familywise error rate & 少なくとも1つの真の帰無仮説を棄却する確率\\
閉検定手順 & closed testing procedure & 共通部分仮説をすべて検定して強い FWER 制御を行う方法\\
用量反応 & dose-response & 用量とともに効果が変化する関係\\
対比 & contrast & 係数の和が0の平均の線形結合\\
混合効果モデル & mixed-effects model & 固定効果と変量効果を含むモデル\\
級内相関係数 & intraclass correlation (ICC) & 同一個体内の測定値の相関\\
事前分布・事後分布 & prior / posterior distribution & Bayes 推測でのパラメータの分布\\
共役事前分布 & conjugate prior & 事後分布が事前分布と同じ分布族になる事前分布\\
信用区間 & credible interval & 事後分布の分位点で作る区間\\
事前予測分布 & prior predictive distribution & データの周辺分布 $\int p(y\mid\theta)p(\theta)\dd\theta$\\
人時間 & person-time & 各人の観察期間の合計．率の分母\\
事後予測分布 & posterior predictive distribution & 将来のデータの予測分布\\
\end{longtable}
\end{center}

% ============================================================
\chapter{過去問対応表}\label[appendix]{app:D}
% ============================================================
必要能力の略号：定＝定義，計＝計算，導＝導出，検＝検定，モ＝モデル化，選＝手法選択，解＝解釈．
問5は4分野共通問題（「共」）で，医薬に関連する部分がある場合のみ対応章を示す．
「F」は設問の一部（または全部）が誘導付きの概念で，\cref{app:F}に一覧したことを表す．

\begin{center}\scriptsize
\setlength{\tabcolsep}{3pt}
\begin{longtable}{@{}ccL{36mm}L{58mm}L{16mm}c@{}}
\toprule
年 & 問 & 主テーマ & 副テーマ & 必要能力 & 章\\
\midrule
\endfirsthead
\toprule
年 & 問 & 主テーマ & 副テーマ & 必要能力 & 章\\
\midrule
\endhead
\bottomrule
\endfoot
2016 & 1 & 対応のある連続データ & 対応のある $t$，符号検定（exact），符号付き順位（正規近似），外れ値の影響 & 検 計 導 解 & \ref{ch:03}\\
2016 & 2 & 4群の比率 & 一様性 $\chi^2_3$，$P=0.077$ の解釈，Cochran--Armitage & 検 導 解 選 & \ref{ch:04}\\
2016 & 3 & 生存時間解析 & KM 曲線，ログランク（$\chi^2=5.62$），Cox 部分尤度 & 計 導 & \ref{ch:09},\ref{ch:10}\\
2016 & 4 & 対数正規と BE & 対数正規の性質（F），クロスオーバーの90\%CI による BE 判定 & 導 計 解 & \ref{ch:13},F\\
2016 & 5共 & 2標本 $t$ 検定 & 群ごとの CI と検定，CI の重なり（約29\%） & 計 検 解 & \ref{ch:01}\\
2017 & 1 & ハザードの理論 & $\lambda=f/S$，$\Lambda=-\log S$，$\Lambda(T)\sim\Ex(1)$，Cox の2重対数表現，乱数生成 & 導 モ & \ref{ch:09},\ref{ch:10}\\
2017 & 2 & 群逐次デザイン & $\alpha$ 消費関数，情報時間，$\Cov[Z(0.4),Z(1)]=\sqrt{0.4}$，限界値 & 計 導 & \ref{ch:14}\\
2017 & 3 & 第II相二段階 & 二段階デザインの $\alpha,\beta$，PET，期待症例数（F） & 導 計 & F\\
2017 & 4 & 層別症例対照 & 曝露 OR，超幾何の平均分散，1段階推定と共通 OR（F），均一性の確認 & 計 導 解 & \ref{ch:04},\ref{ch:07},F\\
2017 & 5共 & 割付法 & 単純無作為化と偏りを抑える割付，CI 幅から $n$ & 計 導 & \ref{ch:02},\ref{ch:12}\\
2018 & 1 & 指数生存モデル & 打ち切り尤度，MLE，Fisher 情報，デルタ法，必要イベント数 & 導 計 & \ref{ch:10},\ref{ch:12}\\
2018 & 2 & マッチドペア2 値 & 多項分布，$\V[\hat\delta]$，CI，制約付き MLE，LRT & 導 検 & \ref{ch:04}\\
2018 & 3 & ROC 解析 & 感度・偽陽性率，経験 ROC，Youden，AUC＝Mann--Whitney & 計 導 解 & \ref{ch:05}\\
2018 & 4 & ロジスティック回帰 & 対数尤度，調整 OR，AIC，Bernoulli の KL，リスクスコアの解釈 & 導 解 & \ref{ch:06},\ref{ch:05}\\
2018 & 5共 & 正規混合 & 平均・分散，二峰性の条件 & 導 & ---\\
2019 & 1 & KM・指数モデル & KM 表，打ち切りのある指数分布の MLE；Nelson--Aalen 型推定・RMST（F） & 計 導 & \ref{ch:09},\ref{ch:10},F\\
2019 & 2 & 交絡と標準化 & 粗 RD，直接標準化・傾向スコア（F） & 計 導 解 & F\\
2019 & 3 & 対応のある診断比較 & 感度・PPV，McNemar 型，8セル多項＋デルタ法 & 計 検 導 & \ref{ch:05},\ref{ch:04}\\
2019 & 4 & 小標本2群比較 & Student $t$ の仮定，順位和の exact 分布，離散性による問題 & 検 計 解 & \ref{ch:03}\\
2019 & 5共 & 適合度と EM & 適合度 $\chi^2$，自由度の根拠，ABO 血液型の EM & 検 導 & \ref{ch:16}\\
2021 & 1 & 競合リスク & $f=\lambda S$，$1-$KM，CIF の導出と推定（F），過大推定の考察 & 導 計 解 & \ref{ch:09},\ref{ch:10},F\\
2021 & 2 & 同等性と BE & 片側 $t$ 検定・IUT（F），90\%CI による BE 判定 & 導 検 解 & \ref{ch:13},F\\
2021 & 3 & 2×2表と層別 & OR と CI，積二項尤度，スコア型 $X$，層別 Cochran 検定，Simpson & 計 導 検 解 & \ref{ch:04},\ref{ch:07}\\
2021 & 4 & 変量効果メタ & 周辺分布，対数尤度，$\hat\mu$，$\V[\hat\mu]$，CI（F） & 導 計 & F\\
2021 & 5共 & 診断の正規モデル & 感度・特異度，PPV・NPV，閾値，有病率の影響 & 計 解 & \ref{ch:05}\\
2022 & 1 & Cox 比例ハザード & 比例性，$S_0^{\exp(\beta x)}$，2重対数プロット（F），HR の CI と解釈 & 導 計 解 & \ref{ch:10},F\\
2022 & 2 & マッチと非可縮性 & 粗 RR・OR，MH 推定量（F），条件付き OR，非可縮性，間接比較 & 計 解 & \ref{ch:04},\ref{ch:07},F\\
2022 & 3 & 線形混合モデル & $\bm V=\bm Z\bm G\bm Z\transp+\bm R$，OLS/GLS，複合対称，被験者内相関 & 導 計 モ & \ref{ch:08}\\
2022 & 4 & Beta--Binomial & 事後分布・モード・平均，周辺分布，事後確率による判定 & 導 計 解 & \ref{ch:17}\\
2022 & 5共 & 平均への回帰 & 2変量正規の条件付き分布（F） & 導 計 解 & F\\
2023 & 1 & ロジスティック回帰 & 対数オッズ，集計データの尤度，MLE，LRT，PPV・NPV & 計 導 検 & \ref{ch:06},\ref{ch:05}\\
2023 & 2 & 単群指数試験の症例数 & MLE，Fisher 情報，Wald 型 $Z$ による症例数（F） & 導 計 & \ref{ch:10},F\\
2023 & 3 & メタアナリシス & 固定効果 MLE，$Q$ の分解と期待値，DL，$I^2$（F） & 導 計 & F\\
2023 & 4 & 反復測定 & $\hat{\bm\mu}$ の分布，共分散構造の AIC，相関行列，変化量の検定 & 計 モ 検 & \ref{ch:08}\\
2023 & 5共 & 条件付き確率と二項検定 & Monty Hall，exact と正規近似 & 計 & ---\\
2024 & 1 & 診断・ROC・McNemar & 感度等の表，台形 AUC，対応のある2 値の分散と McNemar & 計 導 検 & \ref{ch:05},\ref{ch:04}\\
2024 & 2 & ベースライン調整 & 最終値・変化量・ANCOVA の分散（F） & 導 計 解 & F\\
2024 & 3 & 欠測と多重補完 & 欠測と群の $\chi^2$；Rubin のルール・$\E[B]$（F） & 計 導 & \ref{ch:16},F\\
2024 & 4 & 多重性と用量反応 & FWER（F），Bonferroni，分散分析で保護した検定，多対1の相関（F），対比 & 計 導 選 & \ref{ch:14},F\\
2024 & 5共 & 重回帰の抑制 & 正規方程式，決定係数 & 導 & ---\\
2025 & 1 & 順位和検定の検出力 & 位置ずれ，$\E[R_X]$，$\Prob(X>Y)$，検出力の式 & 導 計 & \ref{ch:03}\\
2025 & 2 & ポアソン率と Bayes & 共通率 MLE；率の差の検定・Gamma--Poisson（F） & 導 計 解 & \ref{ch:17},F\\
2025 & 3 & 競合リスクの理論 & 原因別ハザードの関係式（F），原因別 Cox の解釈 & 導 解 & \ref{ch:10},F\\
2025 & 4 & 多重性 & 独立 $p$ 値の FWER・Šidák 型・Simes・相関の影響（F），Bonferroni & 計 導 & \ref{ch:14},F\\
2025 & 5共 & ランダム化回答法 & 誤分類の補正，MLE，$p$ の選択 & 導 解 & ---\\
\end{longtable}
\end{center}

% ============================================================
\chapter{総合チェックリスト}\label[appendix]{app:E}
% ============================================================
試験前に，各項目を\textbf{何も見ずに}説明・導出できるか確認する．できない項目は対応する章に戻る．
\newcommand{\chk}{\item[$\square$]}

\section*{\cref{ch:01,ch:02}　問題の読み方・効果指標}
\begin{itemize}
\chk 設問からデザイン・観測単位・アウトカムの型・推定対象を書き出せる．
\chk $P=0.077$，「CI が1を含む」，「有意差なし」を正しく解釈できる．
\chk 群ごとの CI の重なりと2標本検定の関係（約29\%）を導ける．
\chk $\OR=\RR(1-\pi_0)/(1-\pi_1)$ と，症例対照研究で OR だけが推定できる理由を示せる．
\chk OR の非可縮性を数値例で説明できる．estimand の5要素を挙げられる．
\end{itemize}
\section*{\cref{ch:03}　連続・順位検定}
\begin{itemize}
\chk 符号検定・符号付き順位・順位和（Mann--Whitney）の帰無仮説と仮定を区別できる．
\chk $\mathrm{sign}(Z)$ と $|Z|$ の独立性から $W^+$ の帰無分布を導ける．
\chk 交換可能性と全順位和が一定であることから $\Cov(R_i,R_j)=-\sigma^2/(N-1)$ を導き，$\E[W],\V[W]$ を求められる．
\chk 小標本の exact 分布を数え上げ，最小の P 値を求められる．
\chk $W=U+m(m+1)/2$ と，$\Prob(X>Y)$ を用いた検出力を導ける．
\end{itemize}
\section*{\cref{ch:04}　分割表}
\begin{itemize}
\chk $X^2=N(ad-bc)^2/(n_1n_0m_1m_0)$ をスコア型統計量として導ける．
\chk 二項分布の中心極限定理と多変量デルタ法から $\log\widehat{\RR}$ の漸近正規性を示し，Wald 区間を作れる．
\chk Cochran--Armitage の $\hat\beta$ と $\V[\hat\beta\mid H_0]$ を導き，一様性検定との使い分けを説明できる．
\chk McNemar の分散・検定・LRT（ラグランジュ乗数法）を導ける．条件付き OR $=n_{12}/n_{21}$ を条件付き尤度から導ける．
\chk 超幾何分布の平均・分散から CMH 検定を書ける．
\end{itemize}
\section*{\cref{ch:05}　診断}
\begin{itemize}
\chk 感度・特異度・PPV・NPV を定義し，有病率から PPV を計算できる．
\chk 経験 ROC を描き，台形公式で AUC を計算できる．
\chk AUC $=\Prob(X_1>X_0)$ と経験 AUC＝Mann--Whitney を示せる．
\chk 同一被験者の2検査の感度を McNemar で，PPV を多項＋デルタ法で比較する方法を説明できる．
\end{itemize}
\section*{\cref{ch:06,ch:07}　ロジスティック・交絡}
\begin{itemize}
\chk ロジスティックモデルの対数尤度・スコア・情報行列を導ける．2群なら $2\times2$ 表の対数オッズ比と一致することを示せる．
\chk LRT の自由度，AIC のパラメータ数を正しく数えられる．
\chk 交絡と効果修飾を区別し，Simpson のパラドックスを例示できる．
\chk OR の非可縮性と交絡の違いを説明できる．
\chk Woolf の均一性検定を計算し，Breslow--Day 検定の目的を説明できる．
\end{itemize}
\section*{\cref{ch:08}　縦断・混合モデル}
\begin{itemize}
\chk 反復測定の平均ベクトルの分布から変化量の SE を求められる．
\chk $\bm V=\bm Z\bm G\bm Z\transp+\bm R$，GLS 推定量，変量切片が複合対称を生むことを示せる．
\chk 共分散構造のパラメータ数を数え，AIC を計算できる．
\end{itemize}
\section*{\cref{ch:09,ch:10}　生存時間}
\begin{itemize}
\chk $h=f/S$，$H=-\log S$，$H(T)\sim\Ex(1)$ を証明できる．
\chk KM の表を作り，曲線を描き，中央値を読める．Greenwood の近似公式がどこで近似しているか説明できる．
\chk ログランクの $E_{1j},V_j$ を超幾何分布から導き，分散を時点で足せる理由を説明し，手計算できる．
\chk Cox の部分尤度を書き下し，$S=S_0^{\exp(\beta x)}$ を示せる．
\chk 打ち切りのある指数分布の尤度・MLE・Fisher 情報を導き，デルタ法で HR の検定を作れる．
\chk 原因別ハザード・累積発生関数を定義し，$1-$KM が CIF を過大推定することを示せる．
\chk 原因別 Cox と Fine--Gray の係数の解釈の違いを説明できる．
\end{itemize}
\section*{\cref{ch:12,ch:13,ch:14}　臨床試験デザイン}
\begin{itemize}
\chk $\mu=z_\alpha+z_\beta$ を導き，平均・比率・イベント数の症例数を計算できる．
\chk Schoenfeld の式をログランク分散から導ける．
\chk BE の判定規則（対数変換，GMR の90\%CI $\subset(0.80,1.25)$）を説明し，クロスオーバーデータで判定できる．
\chk $\Corr[Z(t_1),Z(t_2)]=\sqrt{t_1/t_2}$ と $\alpha$ 消費関数による限界値の決め方を説明できる．
\chk FWER，Bonferroni，Holm，閉検定手順を説明でき，分散分析で保護した対比較が FWER を制御する条件を示せる．
\chk 直線的な用量反応に対する最適な対比係数を導ける．
\end{itemize}
\section*{\cref{ch:16,ch:17}　欠測・ベイズ}
\begin{itemize}
\chk MCAR・MAR・MNAR を定義・分類でき，MAR での尤度の無視可能性を示せる．
\chk 完全ケース解析が偏る条件を説明できる．
\chk ABO 血液型の EM の E ステップ・M ステップを導き，試験の公式解答の反復との違いを説明できる．
\chk ポアソン率の MLE と $\V[\log\hat\lambda]\approx1/Y$ を導ける．
\chk Beta--Binomial の事後分布・周辺分布・事後予測分布を導ける．
\chk $\Prob(\pi>p\mid y)=\Prob(\Bin(n+1,p)\le y)$ を用いて事後確率による判定ができる．
\end{itemize}

% ============================================================
\chapter{誘導付きで出題された概念}\label[appendix]{app:F}
% ============================================================
本書では扱わない概念の一覧である．次の概念は，過去問では定義・式・手順が問題文で与えられており，その概念を初めて見る受験者でも誘導に従えば解ける形でしか出題されていない．
そのため本文では扱わない．
いずれも，問題文の式を正しく読み，期待値・分散・尤度の計算を自力で進める力があれば対応できる．
過去問演習では「何が与えられ，何を自分で計算するか」を確認するとよい．

\begin{center}\small
\renewcommand{\arraystretch}{1.15}
\begin{longtable}{@{}L{36mm}L{31mm}L{79mm}@{}}
\toprule
概念 & 出題 & 問題文で与えられていたもの\\
\midrule
\endfirsthead
\toprule
概念 & 出題 & 問題文で与えられていたもの\\
\midrule
\endhead
\bottomrule
\endfoot
\multicolumn{3}{@{}l}{\textbf{連続データ・同等性}}\\
対数正規分布の性質（中央値・密度・平均・分散・CV） & \kakomon{2016}{4}〔1〕〜〔4〕 & 密度の式（「示せ」形式），正規分布の積率母関数\\
TOST・IUT（2つの片側検定） & \kakomon{2021}{2}〔1〕 & 2つの片側仮説，$H_0=H_{01}\cup H_{02}$，「両方棄却で棄却」という手順\\
\multicolumn{3}{@{}l}{\textbf{分割表・交絡}}\\
Peto 型の1段階 log OR，逆分散重み付き共通 OR & \kakomon{2017}{4}〔2〕〜〔4〕 & 超幾何分布の平均・分散，$(O-E)/V$ の形，重み $V$，漸近正規性\\
Mantel--Haenszel の RR・OR 推定量 & \kakomon{2022}{2}〔3〕 & 両推定量の式\\
直接標準化 & \kakomon{2019}{2}〔2〕 & 重み $w_k$ と標準化リスク $R_{xw}$ の定義\\
傾向スコアとバランス特性 & \kakomon{2019}{2}〔3〕〜〔5〕 & 傾向スコアの定義と性質（文章で提示）\\
\multicolumn{3}{@{}l}{\textbf{縦断データ}}\\
最終値・変化量・ANCOVA の比較 & \kakomon{2024}{2} & 2変量正規分布（$\sigma^2$，$\rho$ 既知）と回帰モデル\\
平均への回帰 & \kakomon{2022}{5}（共通） & 2変量正規の条件付き分布の公式\\
Toeplitz 型の共分散構造 & \kakomon{2023}{4}〔2〕 & 共分散行列の形（表で提示）．AIC の式とパラメータ数は前提\\
\multicolumn{3}{@{}l}{\textbf{生存時間}}\\
Nelson--Aalen 型推定 & \kakomon{2019}{1}〔2〕 & $\tilde S(t)=\prod\exp(-d_k/n_k)$ の式\\
境界内平均生存時間（RMST） & \kakomon{2019}{1}〔3〕〜〔5〕 & $X(\tau)=\min(T,\tau)$ の期待値としての定義，指数分布の仮定\\
2重対数プロット & \kakomon{2022}{1}〔3〕 & 軸の定義と「比例ハザードなら平行」という主張\\
原因別ハザード・累積発生関数の関係式 & \kakomon{2021}{1}〔3〕，\kakomon{2025}{3}〔1〕〜〔4〕 & $\lambda_j$，$F_j$，$f_j$ の定義と Weibull 型の原因別ハザードの式（本書は解釈に必要な最小限の定義だけを\cref{sec:10-compete}で扱う）\\
Aalen--Johansen 型の累積発生関数の推定量 & \kakomon{2021}{1}〔4〕 & 推定式 $\hat F_j(t)=\sum d_{ji}/n_{ji}\cdot\hat S(t_{ji}^-)$\\
単群指数試験の症例数（Wald 型 $Z$） & \kakomon{2023}{2}〔4〕〔5〕 & $Z$ の式，$\hat\lambda$ の漸近正規性，情報量を $\lambda_1$ で評価する指示\\
\multicolumn{3}{@{}l}{\textbf{臨床試験デザイン・多重性}}\\
Simon の二段階デザイン & \kakomon{2017}{3} & 手順 $(n_1,r_1,n_2,r)$，記号 $b(x;\pi,n)$，二項確率の表\\
独立な検定の FWER（Šidák 型） & \kakomon{2024}{4}〔1〕，\kakomon{2025}{4}〔1〕 & FWER の定義と独立性（2025年は $p$ 値が独立に一様分布）\\
Dunnett 型の相関，対比どうしの相関 & \kakomon{2024}{4}〔4〕〔6〕 & 各検定統計量 $Z_j$ と対比統計量の式\\
Simes 法 & \kakomon{2025}{4}〔1〕 & 棄却規則 $p_{(1)}\le\alpha/2$ または $p_{(2)}\le\alpha$\\
相関のある検定の FWER & \kakomon{2025}{4}〔2〕 & 2変量正規，$\phi(\rho)$ の定義，証明なしで使ってよい不等式\\
\multicolumn{3}{@{}l}{\textbf{メタアナリシス}}\\
固定効果（逆分散重み），Cochran の $Q$，DerSimonian--Laird，$I^2$ & \kakomon{2023}{3} & 各推定量と $Q$ の定義，$Q$ の分解，「$\E[Q]=Q$ とおく」指示\\
変量効果モデル & \kakomon{2021}{4} & モデル，$\hat\tau^2$，信頼区間の式\\
\multicolumn{3}{@{}l}{\textbf{欠測・ポアソン・ベイズ}}\\
多重補完の Rubin のルール & \kakomon{2024}{3}〔3〕 & $\bar\theta$，$W$，$B$ の定義と $\hat\V[\bar\theta]=W+(1+1/K)B$\\
補完推定量の条件付き分散，$\E[B]$ & \kakomon{2024}{3}〔4〕〔5〕 & 補完の分布（二項分布）の設定\\
ポアソン率の差の $Z$ 検定 & \kakomon{2025}{2}〔2〕 & 正規近似の指定と分散の推定量 $Y/T^2$\\
Gamma--Poisson の事後分布・事後予測分布 & \kakomon{2025}{2}〔3〕〜〔5〕 & ガンマ事前分布の密度（形状・率），パラメータの値，ポアソン性・独立性\\
\end{longtable}
\end{center}
